\documentclass[11pt]{article}
\usepackage[T1]{fontenc}
\usepackage[utf8]{inputenc}
\usepackage{lmodern,amsmath,amssymb,amsthm,microtype}
\usepackage[margin=25mm]{geometry}
\usepackage{hyperref}
\hypersetup{hidelinks,pdftitle={An intrinsic target criterion from a norming atomic reference}}
\setlength{\emergencystretch}{3em}
\newtheorem{theorem}{Theorem}
\newtheorem{lemma}[theorem]{Lemma}
\newtheorem{proposition}[theorem]{Proposition}
\newcommand{\MA}{\operatorname{MA}}
\newcommand{\E}{\mathcal E}
\newcommand{\PSH}{\operatorname{PSH}}
\newcommand{\CPSH}{\operatorname{CPSH}}
\title{An intrinsic target criterion\\from a norming atomic reference}
\author{P16.4 research continuation 220\\Unreviewed partial result}
\date{7 October 2026}
\begin{document}
\maketitle
\begin{abstract}
Assuming the complete independently accepted all-field finite-energy
operator theorem, we prove an intrinsic target-energy criterion when
the continuous reference has a single norming Monge--Ampere atom.
All finite-energy potentials are integrable against a positive target
of the correct raw mass if and only if its continuous-test energy
dual is finite. A convex-series argument gives a uniform unit-energy
bound; the vanishing reference integral then upgrades it to a
square-root estimate and variational coercivity. The hypothesis holds
for standard projective space over every complete characteristic-zero
real-valued non-Archimedean field, including nondiscrete and trivial
valuations. Targets need not be invariant or supported on a skeleton.
No compactness, envelope differentiability or actual equation
attainment is asserted. These new proofs await independent review.
\end{abstract}

\section{Exact prerequisite and norming reference}

Let $K$ be a complete characteristic-zero real-valued
non-Archimedean field, $X/K$ smooth, geometrically integral and
projective of dimension $n\geq1$, and $L$ an actual ample rational
line bundle. Put $T=X^{\mathrm{an}}$ and $V=(L^n)>0$.
Fix a continuous psh reference $\psi$. A scalar $u$ represents the
actual all-point metric $\psi+u$. Its psh definition is the
decreasing-Fubini--Study-net definition in
\cite[Definition 7.2]{BE}; the identically $-\infty$ metric is excluded.
Extended relative energy is the infimum of the continuous psh
majorant energies, with $E(0)=0$, and $\E^1$ consists of actual
potentials of finite extended energy.

Our sole caller mathematical prerequisite is the \emph{complete}
accepted operator manuscript \cite{O}, including its source-preserving
normalization clarifications. We use its actual convex $\E^1$ cone,
energy concavity and decreasing continuity, genuine positive mixed
measures of raw mass $V$, integrability of every $\E^1$ test against
every mixed output, the genuine singular energy identity, affine
curvature calculus, constant law and complete-field pushforward.
No equation, compactness, envelope or comparison result is granted.

We impose the following explicit additional reference hypothesis.
There is a point $\xi\in T$ such that
\begin{equation}\label{eq:anchor}
 \MA(0)=V\delta_\xi,\qquad
 \sup_T u=u(\xi)\quad\hbox{for every scalar FS potential }u.
\end{equation}
These statements concern the reference on this actual bundle.
They are not claimed for a general projective variety or for an
arbitrary atomic measure with several support points.

The norming equality extends to every actual psh potential.
Indeed take an actual FS net $u_i\downarrow u$ and set
$c_i=u_i(\xi)=\sup_Tu_i$. The finite numbers $c_i$ decrease.
If their limit were $-\infty$, the inequalities $u_i\leq c_i$
would force $u$ to be identically $-\infty$. Otherwise their limit
$c$ is finite, $u\leq c$ at every point and $u(\xi)=c$.
Thus $c=\sup_Tu$ and every actual psh potential is finite at $\xi$.

We normalize potentials by $u(\xi)=\sup_Tu=0$, so $u\leq0$.
For such actual $u\in\E^1$ write
\[
                 B(u)=-E(u)\geq0.
\]
If $\mu$ is any positive Radon measure of raw mass $V$, put
\[
 A_\mu(u)=\int_T(-u)\,d\mu\in[0,\infty],\qquad
 J(\mu)=\sup_{v\in\CPSH}\left(E(v)-\int_Tv\,d\mu\right).
\]
Here $\CPSH$ means continuous actual fixed-$L$ potentials. The dual
is unchanged by adding a constant to its test because the two
constant contributions are $cV$. Its value is at least zero.

\begin{theorem}\label{thm:criterion}
Under \eqref{eq:anchor} and the full prerequisite \cite{O}, the
following are equivalent for every positive Radon $\mu$ of mass $V$:
\begin{enumerate}
\item Every actual $u\in\E^1$ is absolutely $\mu$-integrable.
\item The unit-energy integral bound
\[
 U_\mu=\sup\{A_\mu(u):u\in\E^1,\ u(\xi)=0,\ B(u)\leq1\}
\]
is finite.
\item $J(\mu)<\infty$.
\end{enumerate}
If these hold, with $C_n=2^{n+1}$, every normalized $u\in\E^1$
satisfies
\begin{equation}\label{eq:sqrt}
 A_\mu(u)\leq U_\mu\max\{1,\sqrt{C_nB(u)}\}
              \leq U_\mu+U_\mu\sqrt{C_nB(u)}.
\end{equation}
Consequently
\begin{equation}\label{eq:coercive}
 E(u)-\int_Tu\,d\mu
       \leq-B(u)+U_\mu+U_\mu\sqrt{C_nB(u)},\qquad
 J(\mu)\leq U_\mu+\tfrac14 C_nU_\mu^2.
\end{equation}
The functional tends to $-\infty$ whenever the normalized energy
cost tends to infinity, and its supremum over $\E^1$ equals its
continuous-test supremum. None of these assertions says that the
supremum is attained by an actual potential.
\end{theorem}

\section{The quadratic scaling gain}

\begin{lemma}\label{lem:scaling}
For every normalized actual $u\in\E^1$ and $0\leq t\leq1$,
\[
                         B(tu)\leq C_nt^2B(u).
\]
\end{lemma}
\begin{proof}
The convex combination $tu+(1-t)0$ is an actual $\E^1$ potential.
For $0\leq k\leq n$ set
\[
 A_k=\int_T(-u)\,\MA(\underbrace{u,\ldots,u}_{k},
                           \underbrace{0,\ldots,0}_{n-k}).
\]
All numbers are finite and nonnegative by the prerequisite, and
$A_0=-Vu(\xi)=0$ by \eqref{eq:anchor}. The genuine energy identity
gives $B(u)=(n+1)^{-1}\sum_{k=1}^nA_k$.
For $0<t<1$ affine curvature calculus in each slot gives the
\emph{actual} mixed identity
\[
 B(tu)=\frac{t}{n+1}\sum_{j=0}^n\sum_{k=0}^j
       \binom{j}{k}t^k(1-t)^{j-k}A_k.
\]
Its finitely many terms are integrable; no infinite values are
subtracted. The $k=0$ terms vanish. Since $t^{k+1}\leq t^2$ for
$k\geq1$, $0\leq(1-t)^{j-k}\leq1$, and
$\sum_{j=k}^n\binom{j}{k}=\binom{n+1}{k+1}\leq2^{n+1}$,
the displayed expression is at most
$C_nt^2(n+1)^{-1}\sum_{k=1}^nA_k=C_nt^2B(u)$.
The endpoints $t=0,1$ are immediate, and at zero the actual
potential is the reference, without a zero multiple of a pole.
\end{proof}

The absent linear term is the decisive point. Energy concavity
alone would give only $B(tu)\leq tB(u)$ and would not establish the
finite-dual conclusion from a unit-energy bound.

\section{Uniformity from all-test integrability}

Suppose every actual $\E^1$ potential is $\mu$-integrable, but
$U_\mu=\infty$. Choose normalized actual $u_j\in\E^1$ with
$B(u_j)\leq1$ and $A_\mu(u_j)\geq4^j$, for $j\geq1$.
For $N\geq1$ form the actual convex combination
\[
 w_N=\sum_{j=1}^N2^{-j}u_j+2^{-N}0.
\]
Zero coefficients, if any, are omitted. Each $w_N$ belongs to
$\E^1$, $w_N\leq0$, $w_N(\xi)=0$, and $w_{N+1}\leq w_N$.
Concavity gives
\[
                 E(w_N)\geq\sum_{j=1}^N2^{-j}E(u_j)\geq-1.
\]
The decreasing actual psh closure puts the all-point limit $w$
in the psh cone; it is not identically $-\infty$ because
$w(\xi)=0$. Decreasing energy continuity gives $E(w)\geq-1$,
so $w\in\E^1$. Ordinary \emph{sequential} monotone convergence
for the single measure $\mu$ yields
\[
 \int_T(-w)\,d\mu
   =\sum_{j=1}^\infty2^{-j}A_\mu(u_j)
   \geq\sum_{j=1}^\infty2^j=\infty.
\]
This contradicts the hypothesized integrability. Hence $U_\mu$ is
finite. This proof needs no compactness of a family of normalized
potentials and no assertion about arbitrary measurable nets.

Now assume $U_\mu<\infty$. For any normalized $u\in\E^1$, take
$t=1$ if $C_nB(u)\leq1$, and
$t=(C_nB(u))^{-1/2}$ otherwise. In both cases $t>0$ and
Lemma~\ref{lem:scaling} gives $B(tu)\leq1$. Since $tu$ is still
normalized, $tA_\mu(u)=A_\mu(tu)\leq U_\mu$.
This proves \eqref{eq:sqrt}, even when the original integral was
initially allowed to be infinite. It therefore also proves
absolute integrability of every normalized potential. Every other
actual $\E^1$ potential differs by a finite constant, so all-test
integrability follows. For continuous normalized tests
$E(u)-\int u\,d\mu=-B(u)+A_\mu(u)$; maximizing the bound in
\eqref{eq:coercive} over $\sqrt{B(u)}\geq0$ proves finiteness of
$J$ and its asserted upper bound.

\section{Finite dual, singular tests and arbitrary nets}

Conversely let $J(\mu)<\infty$ and take any actual $u\in\E^1$.
Choose continuous psh $u_i\downarrow u$. A tail is bounded above
by a fixed finite constant $M$, because it lies below any one of
its continuous stages. On this tail the continuous nonnegative
functions $M-u_i$ increase to $M-u$. The dual inequality gives
\[
 \int_T(M-u_i)\,d\mu
   \leq MV+J(\mu)-E(u_i)
   \leq MV+J(\mu)-E(u)<\infty.
\]
We use the precise Radon-net integration fact here. The strict
superlevel sets of $M-u_i$ increase as directed open sets and cover
each strict superlevel set of $M-u$. Inner compact approximation
and a finite subcover identify the mass of each limiting superlevel
with the supremum of its stage masses. One common stage approximates
any finite lower layer sum. Truncating the height and refining the
layers therefore shows that the integrals increase to
$\int(M-u)\,d\mu$, including when that integral is infinite.
The uniform finite bound excludes infinity. Since $u$ is bounded
above, this proves absolute $\mu$-integrability.

The same argument, now with the endpoint integral known to be
finite, gives $\int u_i\,d\mu\to\int u\,d\mu$.
The prerequisite supplies $E(u_i)\downarrow E(u)$. Consequently
\[
 E(u)-\int_Tu\,d\mu
       =\lim_i\left(E(u_i)-\int_Tu_i\,d\mu\right)\leq J(\mu).
\]
Continuous tests are included in $\E^1$, proving equality of the
two suprema and completing Theorem~\ref{thm:criterion}.
In particular every genuine output $\mu=\MA(v)$, $v\in\E^1$,
satisfies this target criterion by all-test integrability in
\cite{O}. The converse criterion does not construct an output.

\section{Standard projective space over every complete field}

Take $X=\mathbb P^n_K$, $L=\mathcal O(1)$ and its standard
Fubini--Study reference $\psi$. Let $\xi$ be the Gauss point on
the chart $T_0\ne0$, in coordinates $z_j=T_j/T_0$. For a section
$s$ of $mL$, its expression $s/T_0^m$ is a polynomial of degree
at most $m$. Its value at $\xi$ is the maximum absolute value of
its coefficients. The non-Archimedean triangle inequality in any
projective chart shows that its global $m\psi$-norm is at most
this coefficient maximum; the Gauss point attains equality.

For an actual FS metric defined by sections $s_\alpha$ and real
weights $b_\alpha$, its scalar relative to $\psi$ is
\[
 u=\frac1m\max_\alpha\bigl(\log|s_\alpha|_{m\psi}
                                      +b_\alpha\bigr).
\]
Every section term attains its own maximum at the same $\xi$.
Thus $\sup_Tu=u(\xi)$, including arbitrary real FS weights.
The preceding decreasing-net argument proves the norming identity
for all actual psh endpoints, with literal values at supported points.

The reference measure is $\delta_\xi$, with raw mass $V=1$.
Indeed the standard projective model has reduced irreducible special
fiber and the intersection number of $\mathcal O(1)$ on that fiber
is one. The model formula in \cite[Example 8.11, formula (8.3)]{BE}
therefore has exactly one atom. To apply that formula in its stated
source construction, first pass to a complete algebraically closed
nontrivially valued extension, also when the original valuation is
trivial. The standard model upstairs still has a single reduced
component and its Gauss point maps to $\xi$: Gauss evaluation on
coefficients from $K$ is unchanged. The complete-field pushforward
law gives the claimed original-field measure, without a field-degree
factor. This uses no envelope continuity or density of a value group
in a discretely valued one.

Thus Theorem~\ref{thm:criterion} applies on $\mathbb P^n_K$ over
every complete characteristic-zero real-valued non-Archimedean
field. Its targets are arbitrary positive Radon measures of mass
one subject to the equivalent energy conditions; there is no torus
invariance, skeletal support, carrier metrizability or discreteness
restriction. This explicit domain is a restricted geometric sector,
and is not the entire original arbitrary-variety equation axis.

\section{Remaining equation boundary}

For a general projective variety the reference integral of a
sup-normalized potential need not vanish. The proof above therefore
does not upgrade the general linear scaling bound without an
additional proved reference estimate. Even in the present norming
sector, coercivity alone supplies neither a convergent net of
actual maximizing potentials nor upper semicontinuity in a proved
compactness topology. Envelope orthogonality and both-sign energy
differentiation in arbitrary outer continuous directions, actual
equation attainment and uniqueness of every solution remain to be
established. No original P16.4 closure is claimed. The new
criterion and its complete prerequisite must remain separately
identified when reviewed or reused.

\begin{thebibliography}{9}
\bibitem{O} P16.4 research continuation 217,
\emph{A finite-energy mixed operator over arbitrary complete fields},
complete eight-page manuscript, independently accepted by whole
review 179 on 7 October 2026. Its complete body and harmless
normalization clarifications are the sole caller grant here.
\bibitem{BE} S. Boucksom and D. Eriksson,
\emph{Spaces of norms, determinant of cohomology and Fekete points in
non-Archimedean geometry}, arXiv:1805.01016v3 (2021).
\url{https://arxiv.org/pdf/1805.01016v3}.
\end{thebibliography}
\end{document}
